% year: תשפ"ג
% semester: ב
% moed: א
% source: exams/מבחנים/תשפג/מועד א תשפג סמסטר ב.pdf
% number: 5א
% points: 20
% categories: derivatives, polar-parametric

\begin{question}
מצאו את משוואת המשיק לגרף הפונקציה-
\[ \begin{cases} x=\ln\left(3t^2+e\cdot t-3\right)\\ y=t^3-3t\end{cases} \]
בנקודה $M_0(1,-2)$.
\end{question}

\begin{hint}
מצאו תחילה את ערך הפרמטר $t_0$ המתאים לנקודה $M_0$ (שתי המשוואות צריכות להתקיים בו-זמנית), ואז השתמשו ב-$\frac{dy}{dx}=\frac{y'(t)}{x'(t)}$.
\end{hint}

\begin{solution}
\textbf{מציאת הפרמטר:} מ-$y=-2$: $t^3-3t+2=0\iff(t-1)^2(t+2)=0$, כלומר $t=1$ או $t=-2$.
מ-$x=1$: $3t^2+et-3=e$. עבור $t=1$: $3+e-3=e$ \checkmark. עבור $t=-2$: $12-2e-3=9-2e\neq e$ (כי $e\neq3$). לכן $t_0=1$.

\textbf{הנגזרת:} לפי נוסחת הנגזרת של פונקציה הנתונה פרמטרית,
\[ x'(t)=\frac{6t+e}{3t^2+et-3},\qquad y'(t)=3t^2-3. \]
ב-$t_0=1$: $x'(1)=\frac{6+e}{e}\neq 0$ ו-$y'(1)=0$, ולכן
\[ \frac{dy}{dx}\Big|_{M_0}=\frac{y'(1)}{x'(1)}=0. \]
\textbf{משוואת המשיק:} $y-(-2)=0\cdot(x-1)$, כלומר $\boxed{y=-2}$ (משיק אופקי).
\end{solution}
